\section{Model checking and scorekeeping}

The Bayesian workflow will begin with prior predictive simulation. Priors that
violate stated demographic constraints or assign substantial mass to extreme
paths unsupported by an explicit scenario rationale will be revised before
fitting. Posterior predictive checks will assess whether the model reproduces
salient features of the observed data and whether predictive failures remain
concentrated by region or forecast horizon.

Retrospective evaluation will use rolling information cutoffs. Genuine
vintage-pure hindcasts will use only contemporaneously available inputs;
reconstructed pseudo-out-of-sample exercises will be labelled separately and
will not be conflated with posterior predictive checks. Parameters are
estimated only from information available before the target period, and
predictions are evaluated on later cohort outcomes. Probabilistic comparisons
will use proper scoring rules, including the continuous ranked probability
score already implemented in the project
\citep{gneiting2007}. Calibration and sharpness will be reported separately.

Future predictions will be written as immutable vintages. Each commitment will
identify the quantity, units, resolution source, expected resolution date,
stored predictive distribution, and scoring rule. A new model may supersede an
old one, but it will not overwrite what the old model predicted.
